\section{Results}
We begin with the main workflow-level evidence, then compare selector and admissibility rules directly, and finally add supporting results on scope and mechanism.
The main evidence comes from Camelyon17 ERM, Camelyon17 finetuning, and frozen CivilComments.
ACSIncome is a narrower regression support case, and the remaining analyses are supporting evidence.

\subsection{Core Evidence: Proxy-Selected Reporting Hides Reliability Failures}
Camelyon17 ERM is the cleanest anchor.
Across $10/10$ seeds, proxy-only selection improves the clipped proxy but worsens held-out loss ($+0.055$), ECE ($+0.032$), and the teacher-tail stress diagnostic ($+5.50$).
Mean Brier also worsens ($+0.003$), though only on $8/10$ seeds (Appendix Table~\ref{tab:seed-lure-paired}).
The paired $95\%$ intervals remain fully on the harmful side for loss ($[+0.045,+0.066]$), tail ($[+5.00,+6.04]$), and ECE ($[+0.030,+0.033]$).
A matched epoch-30 check preserves the sign, so the result is not a pure early-stopping artifact.
Figure~\ref{fig:main-divergence} shows the effect visually, Table~\ref{tab:core-cases} summarizes the main deltas, and the appendix gives the seed-matched counts and paired intervals.

\begin{figure*}[t]
  \centering
  \includegraphics[width=\textwidth]{fig_main_divergence_ed}
  \caption{
    Camelyon17 anchor suites under baseline-referenced and within-trajectory selector views.
    Points are mean deltas from the baseline-selected checkpoint.
    Proxy-only, val-loss, and diagnostic-veto selections are evaluated on the suppressive P95 trajectory; the veto may fall back to baseline when no suppressive checkpoint is admissible.
    In both suites, proxy-only moves up on held-out loss; on finetuning, it also moves right on mean held-out accuracy.
    Calibration is a primary reliability check in the tables, while the plotted teacher-tail movement is a fixed-bank stress diagnostic.
    The dashed fixed-horizon arrow tests whether the baseline-referenced harm is only an early-stopping artifact.
  }
  \label{fig:main-divergence}
\end{figure*}

\begin{table}[t]
  \centering
  \small
  \setlength{\tabcolsep}{3.5pt}
  \begin{tabular}{@{}lccccc@{}}
  \toprule
  Case & Hazard & $\Delta$Acc$\uparrow$ & $\Delta$Loss$\downarrow$ & $\Delta$Tail/rel. diag.$\downarrow$ & $\Delta$ECE$\downarrow$ \\
  \midrule
  Camelyon17 ERM & persistent & +0.001 & +0.055 & +5.501 & +0.032 \\
  CelebA ERM & transient & -0.017 & +2.207 & +11.317 & +0.033 \\
  Camelyon17 Finetune & persistent$^{\dagger}$ & \textbf{+0.025} & \textbf{+0.619} & \textbf{+6.226} & \textbf{+0.021} \\
  CivilComments & persistent & -0.004 & +0.104 & +1.829 & +0.070 \\
  ACSIncome & persistent & --- & +0.027 & +0.376 & --- \\
  \bottomrule
\end{tabular}

  \caption{
    Core workflow-hazard cases at P95.
    All deltas are baseline-referenced against the baseline-selected checkpoint from the undistorted run.
    Vision rows use held-out loss and calibration as primary reliability checks; ECE is held-out for the ERM anchors and validation-side for the Camelyon17 finetune anchor, matching the saved calibration artifacts.
    The tail/reliability diagnostic column is a teacher-tail stress diagnostic.
    CivilComments reports official WILDS worst-group loss, and ACSIncome reports tail MSE.
    CelebA ERM is transient; all other rows are persistent, with $\dagger$ marking a qualitative fixed-horizon tail check for Camelyon17 finetuning.
    Baseline held-out losses, for scale, are $0.211$ (Camelyon17 ERM), $1.557$ (CelebA ERM), $0.902$ (Camelyon17 finetune), and $0.258$ (CivilComments).
    Every row also satisfies proxy improvement relative to baseline.
    Appendix Tables~\ref{tab:seed-lure-dominance}--\ref{tab:seed-lure-paired} give seed counts and paired intervals for the classification anchors; the ACSIncome $10/10$ count is backed by the selector artifact.
    Bold marks the strongest ``accuracy up, reliability down'' pattern.
  }
  \label{tab:core-cases}
\end{table}

\subsection{Selector and Veto Checks Change the Reliability Verdict}
Table~\ref{tab:selector-comparison} makes the comparison semantics explicit.
The top panel reports absolute held-out metrics for selectors applied to the two Camelyon17 trajectories.
The bottom panel reports delta-from-baseline selector verdicts for frozen CivilComments and ACSIncome.
The reported-policy column tells the reader where each row's metrics come from: the baseline run, the same suppressive run, a mixed rule that keeps the suppressive checkpoint when admitted and falls back otherwise, or a pure baseline fallback when no suppressive checkpoint is admissible.
The suppressive-admit column records that diagnostic verdict separately.

On the Camelyon anchors, validation-loss-only is safer than proxy-only, but it does not recover the baseline.
Two appendix diagnostics sharpen this point.
First, the Camelyon workflow hazard is visible in standard metrics: Appendix Table~\ref{tab:camelyon-selector-diagnostics} shows harmful held-out loss and reported ECE movement, not only a worse teacher-tail stress diagnostic.
Second, the ERM selector result is stable beyond the main P95 stress test.
On the same suppressive trajectory, proxy-best is worse than validation-loss-best at P95 and P97 on held-out loss, ECE, and the teacher-tail diagnostic; at P99 the contrast attenuates, with loss and Brier favoring proxy-best while ECE and teacher-tail still worsen (Appendix Table~\ref{tab:camelyon-erm-selector-dose}).
Validation-accuracy selection gives the same practical warning, and the leave-one-hospital-out ERM check gives the same standard-metric loss direction across all five held-out folds (Appendix Tables~\ref{tab:camelyon-valacc-selector}--\ref{tab:camelyon-loo-selector-standard}).
The analogous finetune same-trajectory contrast is mixed, so finetuning remains a workflow-level hazard anchor.
The $1.25\times$ veto is therefore only an empirical admissibility diagnostic; when no comparable non-suppressive baseline run is available, the safer operational recommendation is dual reporting of proxy-selected and validation-loss-selected checkpoints.

Camelyon17 finetuning gives the largest mean accuracy-up / reliability-down pattern.
Proxy-only selection raises mean held-out accuracy from $0.683$ to $0.708$, but held-out loss worsens from $0.902$ to $1.520$, the teacher-tail stress diagnostic rises from $7.25$ to $13.48$, and validation-side ECE rises from $0.0308$ to $0.0522$.
The paired intervals remain harmful on loss ($[+0.331,+0.926]$), tail ($[+3.49,+9.07]$), and validation-side ECE ($[+0.018,+0.025]$).
So the main claim in this anchor is a baseline-referenced workflow hazard: reliability degrades under proxy-selected suppressive finetuning, not selector-only causality.
The accuracy increase is descriptive rather than interval-stable, and the fixed-horizon tail check is qualitative only.

On frozen CivilComments, proxy-only and validation-loss-only give the same harmful standard-metric verdict in the bundled summary.
The reversal appears only when a baseline-calibrated validation-loss admissibility rule vetoes the suppressive run and falls back to baseline.
This is the paper's clearest baseline-calibrated run-level admissibility reversal on a fixed suppressive trajectory.
The paired intervals remain harmful on overall loss ($[+0.100,+0.107]$), official worst-group loss ($[+1.785,+1.871]$), and ECE ($[+0.068,+0.071]$).

On ACSIncome, validation-loss-only reduces the winsorization damage ($+0.023$ raw MSE, $+0.282$ tail MSE), but the $1.25\times$ validation-loss veto is slack and coincides with proxy-only.
This is why we treat the veto as a diagnostic rather than a universal selector.
We treat ACSIncome as a narrower regression support case rather than as a co-equal anchor with the two Camelyon suites and frozen CivilComments.
Appendix Table~\ref{tab:guardrail-uncertainty} gives paired intervals for the proxy-only $\rightarrow$ veto reduction on the two Camelyon anchors, and those intervals remain positive for loss and tail reduction at the shown operating point.

\begin{table}[t]
  \centering
  \footnotesize
  \setlength{\tabcolsep}{2.5pt}
  \begin{tabular}{@{}llcccccc@{}}
  \toprule
  Suite & Selector / rule & Reported policy & Supp.\ admit & Sel.\ epoch & Acc / $\Delta$Acc & Loss / $\Delta$Loss & Rel. / $\Delta$Rel \\
  \midrule
  \multicolumn{8}{@{}l}{\emph{Camelyon17 selector comparison (absolute held-out metrics across baseline and suppressive trajectories)}} \\
  \multirow{4}{*}{\shortstack[l]{Camelyon17\\ERM}}
    & Baseline & baseline run & --- & 5.5 & 0.919 & 0.211 & 5.80 \\
    & Proxy-only & suppressive run & 1.00 & 17.0 & 0.920 & 0.266 & 11.31 \\
    & Val-loss & suppressive run & 1.00 & 5.3 & 0.917 & 0.246 & 9.30 \\
    & 1.25x veto & \shortstack[l]{mixed report:\\admit or fallback} & 0.70 & 5.5 & 0.917 & 0.237 & 8.32 \\
  \midrule
  \multirow{4}{*}{\shortstack[l]{Camelyon17\\Finetune}}
    & Baseline & baseline run & --- & 1.6 & 0.683 & 0.902 & 7.25 \\
    & Proxy-only & suppressive run & 1.00 & 3.9 & 0.708 & 1.520 & 13.48 \\
    & Val-loss & suppressive run & 1.00 & 1.8 & 0.677 & 1.511 & 10.05 \\
    & 1.25x veto & \shortstack[l]{mixed report:\\admit or fallback} & 0.40 & 1.4 & 0.693 & 0.970 & 7.71 \\
  \midrule
  \multicolumn{8}{@{}l}{\emph{Cross-modal / regression selector verdicts (deltas from baseline)}} \\
  \multirow{3}{*}{\shortstack[l]{Frozen\\CivilComments}}
    & Proxy-only & suppressive run & 1.00 & --- & $-0.004$ & $+0.104$ & $+1.829$ \\
    & Val-loss & suppressive run & 1.00 & --- & $-0.004$ & $+0.104$ & $+1.829$ \\
    & 1.25x veto & \shortstack[l]{baseline\\fallback} & 0.00 & --- & $+0.000$ & $+0.000$ & $+0.000$ \\
  \midrule
  \multirow{3}{*}{ACSIncome}
    & Proxy-only & suppressive run & 1.00 & --- & --- & $+0.027$ & $+0.376$ \\
    & Val-loss & suppressive run & 1.00 & --- & --- & $+0.023$ & $+0.282$ \\
    & 1.25x veto & suppressive run & 1.00 & --- & --- & $+0.027$ & $+0.376$ \\
  \bottomrule
\end{tabular}

  \caption{
    Selector comparison with semantics made explicit.
    Top panel: absolute held-out metrics for the two Camelyon17 trajectories; bottom panel: delta-from-baseline verdicts for frozen CivilComments and ACSIncome.
    Reported policy identifies whether metrics come from the baseline run, the same suppressive trajectory, an admit-or-fallback mixed report, or pure baseline fallback.
    Suppressive-admit is the diagnostic verdict for the suppressive run; baseline rows are not applicable.
    Frozen CivilComments changes from admitted to vetoed under the validation-loss admissibility diagnostic; on ACSIncome, validation loss is milder but the $1.25\times$ veto is slack.
  }
  \label{tab:selector-comparison}
\end{table}

\subsection{When the Failure Persists and When It Does Not}
CelebA ERM shows that not every hazard is persistent: worst-group loss increases by $+2.21$ and the teacher-tail stress diagnostic by $+11.32$ under proxy-only selection, but fixed epoch-30 deltas reverse sign ($-0.70$ and $-5.93$).
Camelyon17 ERM remains harmful at the fixed horizon, so that case is not just an early-stopping artifact.
Camelyon17 finetuning also keeps harmful fixed-horizon loss, though the fixed-horizon mean loss increase is smaller than at the selected checkpoint ($+0.491$ versus $+0.619$).
Its fixed-horizon tail result should be read only qualitatively because a few seeds explode at epoch~10 (Appendix Table~\ref{tab:persistence-support-app}).
Appendix Table~\ref{tab:persistence-support-app} gives the selected-versus-fixed comparison for every row still labeled persistent, including frozen CivilComments and ACSIncome at its matched final epoch.
Appendix Table~\ref{tab:rw-persistence-bridge} and Figure~\ref{fig:rw-persistence-bridge-app} give the supporting mechanism bridge: CelebA's suppressive orientation relaxes ($\Delta R_w=+0.102$), while Camelyon17's entrenches ($\Delta R_w=-0.091$), so CelebA is a transient counterexample rather than a contradiction.

\subsection{The Failure Is Not Unique to Soft Clipping}
Two ten-seed noisy-label-objective checks on the same Camelyon17 ERM stack address the concern that the result is specific to soft clipping at P95.
Generalized cross-entropy (GCE) with $q=0.7$ gives the strongest corroboration: it is suppressive on this stack at the selected checkpoint ($R_w=0.339$), but it is a standard noisy-label objective rather than a hand-designed clipping stress test.
On the same GCE trajectories, proxy-best selects later checkpoints than validation-loss selection (mean epoch $14.4$ versus $3.3$), lowers the GCE validation proxy by $0.0146$ (95\% CI $[-0.0175,-0.0116]$), and raises held-out accuracy by $+0.0098$ ($[+0.0076,+0.0118]$), yet worsens held-out hospital-2 loss by $+0.0181$ ($[+0.0136,+0.0233]$), ECE by $+0.0129$ ($[+0.0104,+0.0153]$), and teacher-tail by $+3.05$ ($[+2.50,+3.60]$) on $10/10$ seeds (Appendix Table~\ref{tab:gce-q07-extension}).
Baseline-referenced and fixed-epoch-30 GCE contrasts preserve the accuracy-up / reliability-down direction, so this corroboration is not specific to controlled soft clipping.
Hard bootstrapping gives a smaller second noisy-label check: relative to val-loss-only selection inside the same bootstrap runs, proxy-only improves the bootstrap proxy by about $0.010$ but worsens held-out hospital-2 loss by $+0.0109$, teacher-tail by $+1.37$, and ECE by $+0.0110$ (Appendix Table~\ref{tab:bootstrap-camelyon-support}).
A teacher-free held-out stress check gives the same warning: on the top decile of hospital-2 examples fixed by the baseline checkpoint's ordinary cross-entropy loss, proxy-selected SoftClip and GCE increase loss by $+0.94$ and $+0.90$ on $10/10$ seeds (Appendix Table~\ref{tab:baseline-fixed-loss-tail}).
Appendix Table~\ref{tab:camelyon-dose-response-app} adds a same-stack suppression-strength stress test on the two Camelyon anchors.
In the baseline-referenced selected-checkpoint workflow view, Camelyon17 ERM loss / tail harm weakens from P95 to P97 to P99 ($+0.055/+5.50$, $+0.052/+5.16$, $+0.008/+1.76$).
Finetuning remains materially harmful across all three selected regimes.
Together with the same-trajectory P95/P97 diagnostics, P95 is best read as the designated strong stress-test setting rather than as a knife-edge selector artifact.
Appendix Table~\ref{tab:camelyon-erm-selector-dose} separately reports the same-trajectory selector-dose diagnostic, where P95/P97 show clean selector-only harm and P99 attenuates.
This is supporting rather than core evidence because the fixed $1.25\times$ guardrail is slack on this family, but it shows that the selector mismatch is not unique to one clipping rule.

\subsection{The Hazard Concentrates in Suppressive Objectives}
As supporting mechanism evidence, Figure~\ref{fig:mechanism} shows that tail movement is not uniform across objective families on the same Camelyon17 ERM stack: soft clipping has $R_w<1$ and positive tail deltas, while label smoothing and focal loss have $R_w>1$ and negative tail deltas.
On this tested stack, early $R_w$ (epochs 1--5) predicts the sign of the final selected tail delta with $97/100$ seed-level sign accuracy (Appendix Table~\ref{tab:rw-predictive-summary} and Figure~\ref{fig:rw-predictive-summary-app}).
The GCE corroboration ($R_w=0.339$, Appendix Table~\ref{tab:gce-q07-extension}), hard-bootstrap corroboration ($R_w\approx 0.57$, Appendix Table~\ref{tab:bootstrap-camelyon-support}), and end-to-end finetune controls (Appendix Table~\ref{tab:objfam-finetune-support}) support the same suppress-versus-upweight split.
The point is not that any deviation from cross-entropy worsens reliability; it is whether the objective suppresses or amplifies the tail.

The evaluation criterion also changes how objective families rank (Table~\ref{tab:ranking-reversal}).
On frozen Camelyon17 ERM, SoftClip~P95 is the only shown method with positive mean accuracy ($+0.001$), yet it has the worst teacher-tail diagnostic ($+5.50$), a $9.6$ CVaR swing from Focal~4 ($-4.11$).
On end-to-end finetuning, SoftClip has the highest mean accuracy ($+0.047$) but the worst held-out loss ($+0.221$) and teacher-tail diagnostic ($+6.63$), while label smoothing and focal are lower on loss and tail.
Thus accuracy-first and loss/tail-aware readings can favor different methods on the same stack.

\begin{table}[t]
  \centering
  \small
  \begin{tabular}{@{}llrrr@{}}
    \toprule
    Setting & Method & $\Delta$Loss$\downarrow$ & $\Delta$Acc$\uparrow$ & $\Delta$Tail$\downarrow$ \\
    \midrule
    \multirow{3}{*}{\shortstack[l]{Frozen\\ERM}}
      & SoftClip P95    & $+0.055$ & $+0.001$ & $\mathbf{+5.50}$ \\
      & Label Smooth 0.10 & --- & $-0.002$ & $-2.44$ \\
      & Focal 4         & --- & $-0.004$ & $\mathbf{-4.11}$ \\
    \midrule
    \multirow{3}{*}{\shortstack[l]{Separate\\finetune\\control}}
      & SoftClip P95    & $+0.221$ & $\mathbf{+0.047}$ & $+6.63$ \\
      & Label Smooth 0.10 & $-0.276$ & $+0.029$ & $\mathbf{-2.65}$ \\
      & Focal 2         & $\mathbf{-0.280}$ & $-0.026$ & $-2.53$ \\
    \bottomrule
  \end{tabular}
  \caption{
    Within-sweep method-ranking illustration under different reporting criteria.
    Frozen-ERM rows are relative to the ERM baseline; separate finetune-control rows are relative to that sweep's own baseline finetune checkpoint and are for within-table ranking only.
    Focal numbers denote $\gamma$.
    On frozen ERM, the only positive-accuracy row is also worst on teacher-tail; on finetuning, the best-accuracy row is worst on loss and teacher-tail, while label smoothing wins on loss on $10/10$ seeds.
    Bold marks the extreme in each column.
  }
  \label{tab:ranking-reversal}
\end{table}

\begin{figure}[t]
  \centering
  \includegraphics[width=\columnwidth]{fig_mechanism_orientation}
  \caption{
    Same-stack objective-family mechanism on Camelyon17 ERM.
    The vertical boundary at $\log_{10} R_w = 0$ separates tail suppression ($R_w<1$) from tail upweighting ($R_w>1$).
    Focal numbers denote $\gamma$; soft clipping lies on the suppressive side and inflates the tail, while focal loss and label smoothing lie on the upweighting side and improve it.
  }
  \label{fig:mechanism}
\end{figure}
